\section{Limitations}

The framework is a vocabulary and a set of structural relationships, and it is deliberately incomplete. The first limitation is that there is no general algorithm that decides every standing question. The relation depends on an evaluation procedure that may itself be defective, contested or undefined for the transition at hand, and on evidence whose completeness cannot in general be checked from within the procedure that uses it. The relation is a predicate to be assessed, with the arguments that any assessment must supply made explicit, and it is not a function that can be computed once and consulted.

The second limitation concerns what standing cannot remove. It does not eliminate judgment, since choosing a consequence class, a scope or an evaluation procedure is a judgment, and the relation inherits whatever is contestable in those choices. It does not eliminate uncertainty, since a transition can have standing on the available evidence and be mistaken. It does not eliminate institutional conflict, because two bodies with different evidence, procedures and commitments can assign different standing to the same transition without either being mistaken about the facts. A court and a laboratory, or a regulator in one jurisdiction and a regulator in another, can differ in this way, and the relation locates the difference without resolving it.

The third limitation is the regress of evaluators. By the third proposition the use of a procedure does not establish its standing, and a procedure can be assessed by another procedure. It does not follow that every assessment requires an unending chain, and it does not follow that the chain can be closed by proof. In formal systems the chain ends in stipulated rules, and a result is relative to them. In institutions the chain ends in conventions that confer authority on a procedure, and these conventions remain open to contest and revision. In empirical inquiry it ends in practices of replication and error control that are justified partly by their record and partly by the absence of better alternatives. The framework requires that the place where the chain ends be named, and does not claim that naming it settles anything. A document that said otherwise would be offering a foundation it does not possess.

The fourth limitation is that the five transition types are not claimed to be exhaustive. Repair, in which a defect in the record is corrected, revocation, in which standing that was granted is withdrawn, and delegation, in which a body assigns part of its authority to another, are plausible further types. Each interacts with the five given here and with the commitments that accumulate around them, and none is treated, since adding them would require propositions that this document does not have. The first case touches repair informally, and the treatment there should be read as an application of the existing vocabulary and not as a theory of the new type.

The fifth limitation is the relation of the results to the models in which they are proved. The propositions are exhibited by countermodel in small systems, and the programs reproduce them. A countermodel shows that a general claim of preservation, independence or recoverability cannot be made without further assumptions, and it does not show how often the failures it exhibits occur in any real system. The cases apply the vocabulary to real subjects and are as accurate as the vocabulary allows, and none of them claims that the toy models capture those subjects in detail.

The document is for a limited purpose. It makes an existing body of work intelligible as one research program, with one vocabulary and a small set of structural results, and it gives readers a way to see which parts of the corpus are instances of those results, which supply formal machinery that the results only state, and which are separate lines of inquiry. A reader who has read only this document should be able to identify the five transitions, state what standing is relative to, and recognize in a new case which of the five propositions is at issue. What each later essay contributes beyond that is the subject of the accompanying map.
